% Standalone proof note for integration into Part B; AI-produced GPT-6 Astra.
\subsection{The alternating bridge for independently ordered posets}
\begin{proposition}[Proved: finite bridge for induced poset embeddings]
\label{B:poset-bridge}
Let $B,C$ be induced subposets of a finite poset $D$, with an independent
linear order on all structures, and put $F=B\cap C$. Let
$d_0,d_1:D\to E$ be increasing induced-poset embeddings agreeing on $F$.
There is a finite independently ordered poset $Z$, an increasing induced
embedding $E\to Z$, and increasing induced embeddings
$a_0,\ldots,a_{2N}:D\to Z$ with the prescribed endpoints, such that
$a_{2j}|_B=a_{2j+1}|_B$ and
$a_{2j+1}|_C=a_{2j+2}|_C$.
\end{proposition}
\begin{proof}
Use the strict set quotient of one $E$-copy and $2N-1$ copies of $D$
around a circle, identifying alternate overlaps over $B$ and $C$ and
using $d_0,d_1$ at the ends. Write $s=2N$ and require $s>|D|+4$.
As in the explicit fiber calculation in T~A.30, $F$-points have full
visibility and every other point has a proper visibility interval of
at most three consecutive factors. Each factor injects, and adjacent
factor intersections are precisely the prescribed induced copies of
$B$ or $C$. Give the quotient the reflexive transitive closure of all
copied poset relations. We must prove antisymmetry and inducedness.

We use the following elementary path argument, including its reduction.
A path in this generated relation can be described as a succession of
relations inside factors, crossing between consecutive factors at points
of their overlap. If two successive relations were originally recorded
in nonconsecutive factors, insert identity segments along the common
point's visibility interval; for a point outside $F$ this has length at
most two. Split the path at every point of $F$, so every resulting piece
has no $F$-point in its interior. Delete any immediate factor backtrack
$i,j,i$: the endpoints of the segment inside factor $j$ lie in its
induced overlap with factor $i$, so the same relation holds in factor $i$.
This strictly shortens the factor word and does not introduce an $F$
point. Thus a reduced word either has one factor or moves monotonically
along the factor circle.

No such reduced piece can wind around the circle. In fact, a winding
piece crosses the $s-1$ consecutive interior $D$-factors. Inside each
interior factor between two boundary crossings, its boundary labels
belong alternately to $B$ and $C$ and form a comparable pair in $D$.
Equality of these labels would put the boundary point in $F$.
Except possibly at the two endpoints of the piece, the comparison is
therefore strict. All these labels form an increasing chain in the
same finite poset $D$, of length at least $s-3>|D|$, impossible.
The possible three-index visibility through $E$ affects only the
$E$-segment, not this consecutive chain of interior $D$-factors.

Consequently each $F$-free piece misses an overlap. Cut there. Its
entire path lifts to the linear free amalgam of the factors; no extra
identification is crossed by the path. A linear free amalgam of posets
along induced subposets is a poset and contains every factor induced:
for a two-factor amalgam, a cross relation holds exactly when it is
mediated by a base point, and any path returning to one factor contracts
through the induced base; induction proves the linear assertion.

If both endpoints of the original path lie in one factor $T$, and the
path has no $F$-point, choose the endpoint occurrences in $T$ when
forming the factor word (identity transitions within their visibility
intervals suffice). The cut argument proves its relation already in
$T$. If the path does meet $F$, its first and last pieces similarly
give the relations from its first endpoint to a point of $F$, and from
a point of $F$ to its last endpoint, inside $T$. Each intervening piece
has endpoints in $F$ and yields a relation already in $F$, since $F$
is contained induced in every factor of the cut linear amalgam.
Composition proves the entire relation in $T$. Thus every factor is
induced in the generated preorder.

For antisymmetry, a directed cycle disjoint from $F$ reduces to a
closed factor walk. If it does not wind it lifts to a linear free
amalgam and is constant; winding was just excluded. A directed cycle
meeting $F$ splits into pieces whose endpoint relations lie in $F$.
Antisymmetry of $F$ makes all those endpoints equal, and antisymmetry
of each corresponding cut linear amalgam then makes each piece
constant. Hence the quotient preorder is a poset.

The independent linear order is obtained by the interpolation argument
of T~A.30. Embed $E$ increasingly in $\mathbb R$ and interpolate
$p_t=(1-t)d_0+td_1$ on $D$, constant on $F$. For sufficiently large
$N$, alternately hold the $B$-coordinates fixed and move the other
coordinates from $p_{j/N}$ to $p_{(j+1)/N}$; every intermediate map
remains strictly increasing. These maps respect the quotient equations,
so ordering the quotient by their real coordinates and breaking ties
arbitrarily preserves all factor linear orders. There is no requirement
that this independent order extend the poset relation. This gives the
claimed finite bridge.
\end{proof}

\begin{corollary}[Proved: (St) and (Op) for the independently ordered generic poset]
The category of finite posets with an independent linear order and all
increasing induced-poset embeddings satisfies (SA), (FA), (FA$\varnothing$),
(St), and (Op). Its category is total by the all-action theorem in Part B.
\end{corollary}
\begin{proof}
The free-poset amalgam $X\cup_P Z$ has cross relations mediated through
$P$. Its two factors meet exactly in $P$. Given an induced embedding
$u:X\to Y$ over $P$, the map $u\cup\operatorname{id}_Z$ is induced:
a cross comparability $u(x)\leq z$ in the target is equivalent to
$u(x)\leq p\leq z$ for some $p\in P$, which is equivalent to
$x\leq p\leq z$ since $u$ reflects order. Similarly for $z\leq u(x)$.
Choose the $X$-side before the $Z$-side in each gap of the independent
linear order on $P$. This is functorial, and over the empty base it is
disjoint union with concatenated linear order.

For the limit $M$, any finite subsets are strong. Given
$g\in G_{B\cap C}$, take finite $D$ containing $B\cup C$ and finite
$E$ containing $D\cup gD$, and apply the bridge to $1_D,g|_D$.
Embed $Z$ in $M$ over $E$ and extend its $D$-copies by homogeneity.
Consecutive automorphisms differ alternately by elements of $G_B$
and $G_C$, giving $g\in\langle G_B,G_C\rangle$. The other inclusion
is immediate. Minimal finite supports and rigidity of a finite
linear order then give (Op), with the empty support allowed, exactly
as in the proof of Theorem D.2.
\end{proof}
